\en{In this chapter, we prove Clausen's formula (Theorem~\ref{\en{EN}satzclausen}). The proof follows Thomas Clausen's paper \cite[p.~89-91]{\en{EN}clausen} from 1828, but we use modern notations with Pochhammer symbols. Additionally, we prove the hypergeometric differential equations needed for the proof.}%
\de{Ziel dieses Kapitels ist der Beweis der Formel von Clausen (siehe Theorem~\ref{\en{EN}satzclausen}). Der Beweis folgt dem Originalartikel \cite[S.~89-91]{\en{EN}clausen} von Thomas Clausen aus dem Jahr 1828, allerdings in heutiger Notation mit Pochhammer-Symbolen. Auch die für den Beweis benötigten hypergeometrischen Differentialgleichungen werden bewiesen.}

\en{This chapter does not rely on the previous chapters, it is self-contained.}%
\de{Man kann dieses Kapitel unabhängig von den vorherigen Kapiteln lesen.}

\begin{defi}\label{\en{EN}defpoch} \en{For $n\in\mathbb N$ we define the Pochhammer symbol $(a)_n$ as follows:}%
\de{Das Pochhammer-Symbol $(a)_n$ ist für natürliches $n$ wie folgt definiert:}
\begin{align*}
    (a)_0 := 1 \qquad\text{\en{and}\de{und}}\qquad (a)_{n+1} := (a)_n \cdot (a+n)
\end{align*}
\en{This implies $(1)_n=n!$ for all $n$; and $(a)_n = a (a+1) (a+2) \cdots (a+n-1)$, if $n>0$.}%
\de{Hieraus folgt $(1)_n=n!$ für alle $n$; und $(a)_n = a (a+1) (a+2) \cdots (a+n-1)$, falls $n>0$ ist.}
\end{defi}

\begin{defi} \label{\en{EN}defihyp}
\en{The hypergeometric functions ${_2F_1}$ and ${_3F_2}$ are defined as follows:}%
\de{Die hypergeometrische Funktion ${_2F_1}$ und die verallgemeinerte hypergeometrische Funktion ${_3F_2}$ lauten:}
\begin{align*}
    {_2F_1}(a,b;c;z) &= \sum_{n=0}^{\infty} \frac{(a)_n\cdot (b)_n}{(c)_n}\cdot\frac{z^n}{n!}\\
    {_3F_2}(\alpha,\beta,\gamma;\delta,\varepsilon;z) &= \sum_{n=0}^{\infty} \frac{(\alpha)_n\cdot(\beta)_n\cdot(\gamma)_n}{(\delta)_n \cdot (\varepsilon)_n}\cdot\frac{z^n}{n!}
\end{align*}
\end{defi}

\begin{thm}\label{\en{EN}satzkonv}
    ${_2F_1}(a,b;c;z)$ \en{and}\de{und} ${_3F_2}(\alpha,\beta,\gamma;\delta,\varepsilon;z)$ \en{converge absolutely for $|z|<1$.}\de{konvergieren für $|z|<1$ absolut.}
\end{thm}
\begin{proof}
\en{We use the ratio test:}\de{Das folgt aus dem Quotientenkriterium:}
\begin{align*}
    \frac{(a)_{n+1}\cdot (b)_{n+1}}{(c)_{n+1}}\cdot \frac{z^{n+1}}{(n+1)!}\cdot\frac{(c)_n}{(a)_n\cdot (b)_n}\cdot \frac{n!}{z^n} &= \frac{(a+n)(b+n)}{(c+n)(n+1)}\cdot z
\end{align*}
\en{The fraction before $z$ approaches $1$, so the absolute value of the whole expression will be smaller than $1$ for large $n$ if $|z|<1$. The convergence of ${_3F_2}$ is proven in the same way.}%
\de{wobei der Bruch vor $z$ gegen $1$ geht und somit, falls $|z|<1$ ist, der Ausdruck für große $n$ betragsmäßig auch kleiner als $1$ wird. Die Konvergenz von ${_3F_2}$ folgt analog.}
\end{proof}

\begin{thm}\label{\en{EN}koeffvergl}
\en{If}\de{Wenn} $f(z)=\sum_{n=0}^\infty A_n\frac{z^n}{n!}$ \en{is given as a power series, then for all $z$ with absolute convergence of $f$ we get:}\de{als Potenzreihe gegeben ist, dann gilt im absoluten Konvergenzbereich der Reihe:}
\begin{align*}\renewcommand*{\arraystretch}{1.6}
    \begin{array}{rrl|rcl}  f(z)    &:=&\displaystyle\sum_{n=0}^\infty A_n\frac{z^n}{n!}\\
                    zf'(z)  &=& \displaystyle\sum_{n=0}^\infty n A_n\frac{z^n}{n!}
                    &    f'(z) &=&\displaystyle\sum_{n=0}^\infty A_{n+1}\frac{z^n}{n!}\\
                    z^2f''(z) &=& \displaystyle\sum_{n=0}^\infty n(n-1) A_n\frac{z^n}{n!}
                    & zf''(z) &=&\displaystyle\sum_{n=0}^\infty n A_{n+1}\frac{z^n}{n!}\\
                    z^3f'''(z) &=& \displaystyle\sum_{n=0}^\infty n(n-1)(n-2) A_n\frac{z^n}{n!}~~~
                    &~~ z^2f'''(z) &=&\displaystyle\sum_{n=0}^\infty n(n-1) A_{n+1}\frac{z^n}{n!}
    \end{array}
\end{align*}
\end{thm}
\begin{proof}
  \en{Since the power series $f(z)$ converges absolutely according to our premises, we may interchange summation and derivation.
  The identities are a direct consequence of this and the definition of $f(z)$:
  
  The reason we didn't reduce the fractions in the left identities with $n$ etc.~is that it will make equating coefficients easier.
  For the identites on the right side, we skip the first summand in $f'(z)=\sum_{n=0}^\infty A_n\frac{n\cdot z^{n-1}}{n!}=\sum_{n=1}^\infty A_n\frac{n\cdot z^{n-1}}{n!}=\sum_{n=1}^\infty A_n\frac{z^{n-1}}{(n-1)!}$ (which is zero anyway), then we reduce the fraction by $n$.
  Then we do an index shift by $1$ and get the first identity on the right side, i.e.~$f'(z)=\sum_{m=0}^\infty A_{m+1}\frac{z^m}{m!}$.
  For the further formulae on the right side, we don't use any further index shifts.\pagebreak
  }%
  \de{Weil die Reihe $f(z)$ nach Voraussetzung absolut konvergiert, dürfen wir Summation und Ableitung vertauschen.
  Die Formeln folgen dann direkt aus der Definition von $f(z)$.
  Bei den linken Formeln wird nicht gekürzt, damit man es später beim Koeffizientenvergleich leichter hat.
  Für die Formeln der rechten Spalte wird zunächst bei $f'(z)=\sum_{n=0}^\infty A_n\frac{n\cdot z^{n-1}}{n!}=\sum_{n=1}^\infty A_n\frac{n\cdot z^{n-1}}{n!}=\sum_{n=1}^\infty A_n\frac{z^{n-1}}{(n-1)!}$ der nullte Summand weggelassen (der sowieso Null ist), dann wird mit $n$ gekürzt.
  Jetzt folgt noch eine Indexverschiebung um $1$ und wir erhalten die oberste Formel der rechten Spalte, nämlich $f'(z)=\sum_{m=0}^\infty A_{m+1}\frac{z^m}{m!}$.
  Für die weiteren Formeln der rechten Spalte wird wieder aufs Kürzen und Indexverschieben verzichtet.}
  \end{proof}


\pagebreak

\begin{theo}\label{\en{EN}dgl2f1}
\en{The hypergeometric function $f(z)={_2F_1}(a,b;c;z)$ satisfies the "hypergeometric differential equation":}%
\de{Die hypergeometrische Funktion $f(z)={_2F_1}(a,b;c;z)$ erfüllt die hypergeometrische Differentialgleichung:}
$$z(z-1) f''(z) + \left[(a+b+1) z - c\right] f'(z) + ab f(z) = 0$$
\end{theo}
\begin{proof}
    \en{We prove this by equating coefficients.
    The hypergeometric series can be written as $\displaystyle f(z)=\sum_{n=0}^\infty A_n\frac{z^n}{n!}$ with $\displaystyle A_n:=\frac{(a)_n\cdot(b)_n}{(c)_n}$.
    The Definition~\ref{\en{EN}defpoch} of the Pochhammer symbols yields $(a)_{n+1}=(a)_n\cdot(a+n)$ and thus }%
    \de{Wir führen den Beweis mit Hilfe eines Koeffizientenvergleichs.
    Die hypergeometrische Funktion lautet $\displaystyle f(z)=\sum_{n=0}^\infty A_n\frac{z^n}{n!}$ mit $\displaystyle A_n:=\frac{(a)_n\cdot(b)_n}{(c)_n}$.
    Die Definition~\ref{\en{EN}defpoch} der Pochhammersymbole sagt $(a)_{n+1}=(a)_n\cdot(a+n)$ und somit }
    \begin{align*}
        A_{n+1}&=\frac{(a+n)(b+n)}{(c+n)}\cdot A_n\\
        \Longrightarrow (c+n)\cdot A_{n+1} &= (n^2+(a+b)n+ab)\cdot A_n\\
        \Longrightarrow (c+n)\cdot A_{n+1} &= (n(n-1)+(a+b+1)n+ab)\cdot A_n\\
        \Longrightarrow c\cdot A_{n+1} + n\cdot A_{n+1} &= n(n-1)\cdot A_n+(a+b+1) n \cdot A_n+ab\cdot A_n
    \end{align*}
    \en{Now we recognize the coefficients of Prop.~\ref{\en{EN}koeffvergl} and get:}%
    \de{Wenn wir nun Satz~\ref{\en{EN}koeffvergl} verwenden, erkennen wir an dieser Koeffizientengleichung:}
    \begin{align*}
        cf'(z) + zf''(z) = z^2f''(z) + (a+b+1)zf'(z)+abf(z)\\
        \Longrightarrow\qquad z(z-1) f''(z) + \left[(a+b+1) z - c\right] f'(z) + ab f(z) = 0
    \end{align*}
    \en{Thus we have proven that ${_2F_1}$ satisfies the hypergeometric differential equation.}%
    \de{Somit ist bewiesen, dass die hypergeometrische Funktion die genannte Differentialgleichung erfüllt.}
\end{proof}

\begin{thm}\label{\en{EN}dgl3f2}
\en{The hypergeometric function $g(z)={_3F_2}(\alpha,\beta,\gamma;\delta,\varepsilon;z)$ satisfies this differential equation:}%
\de{Die verallgemeinerte hypergeometrische Funktion $g(z)={_3F_2}(\alpha,\beta,\gamma;\linebreak[3]\delta,\varepsilon;z)$ erfüllt die folgende Differentialgleichung:}
\begin{align*}
    (z^3-z^2)\cdot g'''(z) + [(\alpha+\beta+\gamma+3)z^2-(\delta+\varepsilon+1)z]\cdot g''(z)&\\
    +~[(1+\alpha+\beta+\gamma+\alpha\beta+\alpha\gamma+\beta\gamma)z-\delta\varepsilon]\cdot g'(z)+\alpha\beta\gamma\cdot g(z)&=0
\end{align*}
\end{thm}
\begin{proof}
    \en{We prove this like in Thm.~\ref{\en{EN}dgl2f1} by equating coefficients.
    The hypergeometric function now is $\displaystyle g(z)=\sum_{n=0}^\infty A_n\frac{z^n}{n!}$ with the coefficients $\displaystyle A_n:=\frac{(\alpha)_n\cdot(\beta)_n\cdot(\gamma)_n}{(\delta)_n \cdot (\varepsilon)_n}$.
    The Definition~\ref{\en{EN}defpoch} of the Pochhammer symbols yields $(a)_{n+1}=(a)_n\cdot(a+n)$ and thus}%
    \de{Wir verwenden wieder einen Koeffizientenvergleich als Beweis, völlig analog zum Beweis von Thm.~\ref{\en{EN}dgl2f1}.
    Die verallgemeinerte hypergeometrische Funktion lautet $\displaystyle g(z)=\sum_{n=0}^\infty A_n\frac{z^n}{n!}$ mit den Koeffizienten $\displaystyle A_n:=\frac{(\alpha)_n\cdot(\beta)_n\cdot(\gamma)_n}{(\delta)_n \cdot (\varepsilon)_n}$.
    Die Definition~\ref{\en{EN}defpoch} der Pochhammersymbole sagt $(a)_{n+1}=(a)_n\cdot(a+n)$ und somit }
    \begin{align*}
        A_{n+1}&=\frac{(\alpha+n)(\beta+n)(\gamma+n)}{(\delta+n)(\varepsilon+n)}\cdot A_n\\
        \Longrightarrow (\delta+n)(\varepsilon+n)\cdot A_{n+1} &= (\alpha+n)(\beta+n)(\gamma+n)\cdot A_n\\
        \Longrightarrow [n^2+(\delta+\varepsilon)n+\delta\varepsilon] A_{n+1} &= [n^3+(\alpha+\beta+\gamma)n^2+(\alpha\beta+\alpha\gamma+\beta\gamma)n+\alpha\beta\gamma] A_n
    \end{align*}
    \en{But it is $n^2=n(n-1)+1n$ and $n^3=n(n-1)(n-2)+3n^2-2n$, so we get}%
    \de{Jetzt ist aber $n^2=n(n-1)+1n$ und $n^3=n(n-1)(n-2)+3n^2-2n$, also gilt}
    \begin{align*}
         &~[n(n-1)+(\delta+\varepsilon+1)n+\delta\varepsilon]\cdot  A_{n+1} \\
        =&~[n(n-1)(n-2)+(\alpha+\beta+\gamma+3)n^2+(\alpha\beta+\alpha\gamma+\beta\gamma-2)n+\alpha\beta\gamma]\cdot  A_n\\
        \Longrightarrow~~~~&~n(n-1)  A_{n+1}+(\delta+\varepsilon+1)n  A_{n+1}+\delta\varepsilon  A_{n+1} = n(n-1)(n-2) A_n\\
        &\qquad\qquad+~(\alpha+\beta+\gamma+3)n^2 A_n+(\alpha\beta+\alpha\gamma+\beta\gamma-2)n A_n+\alpha\beta\gamma A_n\\
        \Longrightarrow~~~~&~n(n-1)  A_{n+1}+(\delta+\varepsilon+1)n  A_{n+1}+\delta\varepsilon  A_{n+1} = n(n-1)(n-2) A_n\\
        &\qquad\qquad+~(\alpha+\beta+\gamma+3)n(n-1) A_n\\
        &\qquad\qquad+~(\alpha\beta+\alpha\gamma+\beta\gamma-2+\alpha+\beta+\gamma+3)n A_n+\alpha\beta\gamma A_n
    \end{align*}
    \en{Again we recognize the coefficients of Prop.~\ref{\en{EN}koeffvergl} and get:}%
    \de{Wenn wir nun Satz~\ref{\en{EN}koeffvergl} verwenden, erkennen wir an dieser Koeffizientengleichung:}
    \begin{align*}
        &z^2\cdot g'''(z)+(\delta+\varepsilon+1)z\cdot g''(z)+\delta\varepsilon\cdot g'(z)\\
        &\qquad= z^3\cdot g'''(z)+(\alpha+\beta+\gamma+3)z^2g''(z)\\
        &\qquad+(\alpha\beta+\alpha\gamma+\beta\gamma+\alpha+\beta+\gamma+1)z\cdot g'(z)+\alpha\beta\gamma\cdot g(z)\\
        \Longrightarrow~~~~&~\left[z^3-z^2\right]\cdot g'''(z) + [(\alpha+\beta+\gamma+3)z^2-(\delta+\varepsilon+1)z]\cdot g''(z)\\
        &\qquad+[(\alpha\beta+\alpha\gamma+\beta\gamma+\alpha+\beta+\gamma+1)z-\delta\varepsilon]\cdot g'(z)+\alpha\beta\gamma\cdot g(z)=0
    \end{align*}
    \en{Thus we have proven that ${_3F_2}$ satisfies said differential equation.}%
    \de{Somit ist bewiesen, dass die verallgemeinerte hypergeometrische Funktion die genannte Differentialgleichung erfüllt.}
\end{proof}


\begin{theo}\label{\en{EN}satzclausen}
\en{The following formula (published and proven in 1828 by Thomas Clausen) applies:}\de{Es gilt die Formel von Thomas Clausen aus dem Jahr 1828, nämlich}
$$\left({_2F_1}\lk a,b;a+b+\frac 1 2;z\rk\right)^2 = {_3F_2}\lk 2a,2b,a+b;2a+2b,a+b+\frac 1 2;z\rk.$$
\en{Setting $a=\frac{1}{12}$ and $b=\frac{5}{12}$ yields:}\de{Insbesondere gilt für $a=\frac{1}{12}$ und $b=\frac{5}{12}$:}
$$\left({_2F_1}\lk\frac{1}{12},\frac{5}{12};1;z\rk\right)^2 = {_3F_2}\lk\frac{1}{6},\frac{5}{6},\frac{1}{2};1,1;z\rk.$$
\end{theo}
\begin{proof}
    \en{We prove that both sides of this equation satisfy the same third order differential equation:}%
    \de{Wir zeigen zunächst, dass beide Seiten der Gleichung derselben Differentialgleichung dritter Ordnung genügen.}% (vgl.~\cite{\en{EN}andrews_askey_roy_1999}, Übungsaufgabe 2.13, S.~116).
    
    \en{First we look at the right side, which we call $g(z)$. Here we recognize the hypergeometric function ${_3F_2}$ with its differential equation from Prop.~\ref{\en{EN}dgl3f2}. Setting $\alpha=2a$, $\beta=2b$, $\gamma=a+b$, $\delta=2a+2b$ and $\varepsilon=a+b+\frac 1 2$ we get:}%
    \de{Zunächst betrachten wir die rechte Seite, die wir $g(z)$ nennen. Das ist die verallgemeinerte hypergeometrische Funktion ${_3F_2}$, für die wir in Satz~\ref{\en{EN}dgl3f2} bereits eine Differentialgleichung bewiesen haben. Wir müssen nur noch $\alpha=2a$, $\beta=2b$, $\gamma=a+b$, $\delta=2a+2b$ und $\varepsilon=a+b+\frac 1 2$ einsetzen und erhalten:}
    \begin{align}
        \left[z^3-z^2\right]\cdot g'''(z) + \left[3\left(a+b+1\right)z^2-3\left(a+b+\frac 1 2\right)z\right]\cdot g''(z)&\nonumber\\
        +~\left[\left(1+3a+3b+8ab+2a^2+2b^2\right)z-\left(a+b\right)\left(2a+2b+1\right)\right]\cdot g'(z)&\label{\en{EN}ue13}\\
        +~4ab\left(a+b\right)\cdot g(z)&=0\nonumber
    \end{align}
    
    \en{Next we look at the left side, which we call $h(z):=\left({_2F_1}\lk a,b;a+b+\frac 1 2;z\rk\right)^2$. Here it gets more complicated, since we have the square of a power series. But we will show that $h(z)$ also satisfies the differential equation~(\ref{\en{EN}ue13}).}%
    \de{Jetzt kommen wir zur linken Seite, die wir $h(z)$ nennen. Hier steht das Quadrat einer Potenzreihe, deshalb wird es wesentlich komplizierter. Wir werden zeigen, dass auch $h(z):=\left({_2F_1}\lk a,b;a+b+\frac 1 2;z\rk\right)^2$ eine Lösung der Differentialgleichung~(\ref{\en{EN}ue13}) ist.}
    
    \en{We start with $f(z)={_2F_1}\left(a,b;a+b+\frac 1 2;z\right)$, for which we have the differential equation from Thm.~\ref{\en{EN}dgl2f1} with $c=a+b+\frac 1 2$:}%
    \de{Für $f(z)={_2F_1}\left(a,b;a+b+\frac 1 2;z\right)$ gilt die Differentialgleichung aus Thm.~\ref{\en{EN}dgl2f1}, wobei $c=a+b+\frac 1 2$ ist:}
    \begin{align}
        \left(z^2-z\right) f''(z) + \left[(a+b+1) z - c\right] f'(z) + ab f(z) &= 0\qquad\left|~\cdot~z\right.\label{\en{EN}C0}\\
        \left(z^3-z^2\right) f''(z) + \left[(a+b+1) z^2 - cz\right] f'(z) + abz f(z) &= 0\qquad\left|~\frac{d}{dz}\right.\label{\en{EN}C1}\\
        \left(z^3-z^2\right) f'''(z) + \left[(a+b+4) z^2 - \left(c+2\right)z\right] f''(z)&\nonumber\\
        +~\left[\left(ab+2a+2b+2\right)z-c\right]f'(z) + ab f(z) &= 0\label{\en{EN}C2}
    \end{align}
    \en{In the last step we derived the equation by $z$ minding the product rule and grouped similar terms.
    Now we form a linear combination of the equations~(\ref{\en{EN}C0}),~(\ref{\en{EN}C1}) and~(\ref{\en{EN}C2}), as Clausen suggested in \cite{\en{EN}clausen}:}%
    \de{Im letzten Schritt haben wir die Gleichung unter Beachtung der Produktregel nach $z$ abgeleitet und dann ähnliche Terme gruppiert.
    Nun bilden wir (wie von Clausen in \cite{\en{EN}clausen} vorgeschlagen) eine Linearkombination der Gleichungen~(\ref{\en{EN}C0}),~(\ref{\en{EN}C1}) und~(\ref{\en{EN}C2}), nämlich:}
    $$(2a+2b-1)\cdot2f(z)\cdot(\ref{\en{EN}C0}) + 6f'(z)\cdot(\ref{\en{EN}C1})+ 2f(z)\cdot(\ref{\en{EN}C2})$$


    \pagebreak
    \en{This linear combination reads:\\}\de{Diese Linearkombination lautet ausgeschrieben:\\}
    \noindent\scalebox{0.93}{\begin{minipage}{\the\textwidth}
\begin{align}
    0 = (2a+2b-1)\cdot 2f \cdot \left[ \underline{\underline{(z^2-z) f''}} + \underline{\left((a+b+1)z-c\right)f'} + abf\right]&\nonumber\\
    +~6f'\cdot\left[\underline{\underline{\underline{(z^3-z^2) f''}}} + \underline{\underline{\left((a+b+1)z^2-cz\right)f'}} + \underline{abz f}\right]&\label{\en{EN}unterstr}\\
    +~2f\cdot\left[\underline{\underline{\underline{(z^3-z^2) f'''}}} + \underline{\underline{\left((a+b+4)z^2-(c+2)z\right)f''}} + \underline{\left((ab+2a+2b+2)z-c\right)f'} + ab f\right] &\nonumber
\end{align}
\end{minipage}}\\[1ex]

\en{Now we will combine the different terms in equation~(\ref{\en{EN}unterstr}) in order to get equation~(\ref{\en{EN}linkombclausen}) later on.
The terms that are three times underlined in~(\ref{\en{EN}unterstr}) contain the third derivatives:}%
\de{Jetzt werden wir die verschiedenen Terme in Gleichung~(\ref{\en{EN}unterstr}) zusammenfassen, um weiter unten Gleichung~(\ref{\en{EN}linkombclausen}) zu erhalten.
Die dreifach unterstrichenen Summanden in~(\ref{\en{EN}unterstr}) enthalten die dritten Ableitungen:}
$$\underline{\underline{\underline{\left\{2f f''' + 6 f' f''\right\}}}}\cdot\left[z^3-z^2\right]$$
\en{The terms in~(\ref{\en{EN}unterstr}) that are twice underlined contain the second derivatives:}%
\de{Die zweifach unterstrichenen Summanden in~(\ref{\en{EN}unterstr}) enthalten die zweiten Ableitungen:}
\begin{align*}
    &\underline{\underline{\left\{2f f''\right\}}}\cdot \underbrace{\left[(2a+2b-1)\cdot(z^2-z)+(a+b+4)z^2-(c+2)z\right]}_{=:A_1}\\
    +~&\underline{\underline{\left\{2f'^2\right\}}}\cdot \underbrace{\left[3(a+b+1)z^2-3(a+b+1/2)z\right]}_{=:A_2}\\
    =~&\underline{\underline{\left\{2f'f'' + 2f'^2\right\}}}\cdot \left[3(a+b+1)z^2-3(a+b+1/2)z\right],\\
    \text{\en{since}\de{denn} }A_1 =~ & (2a+2b-1+a+b+4)z^2-(2a+2b-1+a+b+1/2+2)z = A_2
\end{align*}
\en{The terms that are underlined once contain the first derivatives:}%
\de{Die einfach unterstrichenen Summanden in~(\ref{\en{EN}unterstr}) enthalten die ersten Ableitungen:}
\begin{align*}
&\underline{\left\{2f f'\right\}}\cdot \left[(2a+2b-1)\cdot\left((a+b+1)z-c\right)+3abz+\left((ab+2a+2b+2)z-c\right)\right]\\
=~&\underline{\left\{2f f'\right\}}\cdot \left[ \left( (2a+2b-1)(a+b+1)+4ab+2a+2b+2 \right)z-\left((2a+2b-1)c+c\right) \right]\\
=~&\underline{\left\{2f f'\right\}}\cdot \left[ \left( 1+3a+3b+8ab+2a^2+2b^2 \right)z-(a+b)(2a+2b+1) \right]
\end{align*}
\en{Finally, the terms that aren't underlined contain no derivatives:}%
\de{Die nicht unterstrichenen Summanden in~(\ref{\en{EN}unterstr}) enthalten keine Ableitungen:}
\begin{align*}
\left\{f^2\right\}\cdot \left[2ab\cdot(2a+2b-1)+2ab\right] = \left\{f^2\right\}\cdot \left[4ab(a+b)\right]
\end{align*}
    \en{If we combine all this, equation~(\ref{\en{EN}unterstr}) yields:}%
    \de{Insgesamt geht Gleichung~(\ref{\en{EN}unterstr}) also über in:}
    \begin{align}
        &\left[z^3-z^2\right]\cdot\left\{2ff'''+6f'f''\right\}\nonumber\\
        +&\left[3\left(a+b+1\right)z^2-3\left(a+b+\frac 1 2\right)z\right]\cdot \left\{2ff''+2f'^2\right\}\label{\en{EN}linkombclausen}\\
        +&\left[\left(1+3a+3b+8ab+2a^2+2b^2\right)z-\left(a+b\right)\left(2a+2b+1\right)\right]\cdot \left\{2ff'\right\}\nonumber\\
        +&\left[4ab\left(a+b\right)\right]\cdot \left\{f^2\right\} =0\nonumber
    \end{align}
    \en{and here we recognize the content of the square brackets from equation~(\ref{\en{EN}ue13}).\\
    Next we use $h(z)=\left(f(z)\right)^2$:}%
    \de{und wir erkennen, dass der Inhalt der eckigen Klammern genau so auch in Gleichung~(\ref{\en{EN}ue13}) vorkam.
    Nun gilt aber auch noch, dass $h(z)=\left(f(z)\right)^2$, und somit:}
    \begin{align}
        h(z) &= (f(z))^2\nonumber\\
        h'(z) &= 2f(z)f'(z)\nonumber\\
        h''(z) &= 2f(z)f''(z)+2(f'(z))^2\label{\en{EN}ablH}\\
        h'''(z) &= 2f'(z)f''(z)+2f(z)f'''(z)+4f'(z)f''(z)\nonumber\\
                &= 2f(z)f'''(z)+6f'(z)f''(z)\nonumber
    \end{align}
    \en{If we now replace the curly brackets in~(\ref{\en{EN}linkombclausen}) by these derivatives of $h(z)$, we realize that $h(z)=\left({_2F_1}\lk a,b;a+b+\frac 1 2;z\rk\right)^2$ is another solution to the differential equation~(\ref{\en{EN}ue13}).\pagebreak}%
    \de{Wir können also für die geschweiften Klammern in~(\ref{\en{EN}linkombclausen}) die Ableitungen von $h(z)$ einsetzen und haben bewiesen, dass auch $h(z)=\left({_2F_1}\lk a,b;a+b+\frac 1 2;z\rk\right)^2$ eine Lösung der Differentialgleichung~(\ref{\en{EN}ue13}) ist.}
     
   
     
    \en{We still have to show that these two solutions are the same, that we have $h(z)=g(z)$.
    Since they both are solutions to the same ordinary differential equation of third order, it suffices to show that $h(0)=g(0)$ and $h'(0)=g'(0)$ and $h''(0)=g''(0)$.
    In order to prove this, we use the notation $A_n:=\frac{(a)_n\cdot(b)_n}{(a+b+1/2)_n}$ for the coefficients of $f(z)$.
    The first three of these coefficients are (cf.~Def.~\ref{\en{EN}defpoch}):}%
    \de{Wir müssen jetzt noch zeigen, dass die beiden Lösungen sogar gleich sind, also dass $h(z) = g(z)$ gilt.
    Da es sich um eine Differentialgleichung dritter Ordnung handelt, müssen wir z.B.~für $z=0$ zeigen, dass die Funktionswerte und die ersten beiden Ableitungen paarweise übereinstimmen.
    Hierfür verwenden wir die Notation $A_n:=\frac{(a)_n\cdot(b)_n}{(a+b+1/2)_n}$ für die Koeffizienten von $f(z)$.
    Die ersten drei dieser Koeffizienten lauten (vgl.~Def.~\ref{\en{EN}defpoch}):}
    \begin{align*}
        A_0 &= \frac{(a)_0\cdot(b)_0}{(a+b+1/2)_0}=\frac{1\cdot 1} 1 = 1\\
        A_1 &= \frac{(a)_1\cdot(b)_1}{(a+b+1/2)_1}=\frac{a\cdot b} {a+b+\frac 1 2}\\
        A_2 &= \frac{(a)_2\cdot(b)_2}{(a+b+1/2)_2}=\frac{a(a+1)\cdot b(b+1)} {\left(a+b+\frac 1 2\right)\left(a+b+\frac 3 2\right)}
    \end{align*}
    \en{Then we use the derivatives in~(\ref{\en{EN}ablH}) and obtain the following values of $h(z)=(f(z))^2$:}%
    \de{Somit gilt für $h(z)=(f(z))^2$ wegen der Ableitungsdarstellungen~(\ref{\en{EN}ablH}):}
    \begin{align*}
        h(0) &= (f(0))^2 = A_0^2 = 1\\
        h'(0) &= 2 f(0) f'(0) = 2 A_0 A_1 = \frac{2 ab} {a+b+\frac 1 2}\\
        h''(0) &= 2f(0)f''(0) + 2(f'(0))^2 = 2A_0A_2+2A_1^2\\
            &=\frac{2ab(a+1)(b+1)} {(a+b+1/2)(a+b+3/2)} + \frac{2a^2b^2} {(a+b+1/2)^2}\\
            &=\frac{ab(4a^2b+4ab^2+8ab+2a^2+2b^2+3a+3b+1)}{(a+b+1/2)^2(a+b+3/2)}
    \end{align*}
    \en{For the other solution $g(z)={_3F_2}(\alpha,\beta,\gamma;\delta,\varepsilon;z)$ it holds:}%
    \de{Für die andere Lösung $g(z)={_3F_2}(\alpha,\beta,\gamma;\delta,\varepsilon;z)$ gilt:}
    \begin{align*}
        g(0)&=\frac{(\alpha)_0\cdot(\beta)_0\cdot(\gamma)_0}{(\delta)_0 \cdot (\varepsilon)_0} = 1 = h(0)\\
        g'(0)&=\frac{(\alpha)_1\cdot(\beta)_1\cdot(\gamma)_1}{(\delta)_1 \cdot (\varepsilon)_1} = \frac{\alpha\beta\gamma}{\delta\varepsilon}=\frac{2a\cdot2b\cdot(a+b)}{(2a+2b)\cdot\left(a+b+\frac 1 2\right)}= \frac{2 ab} {a+b+\frac 1 2}=h'(0)\\
        g''(0)&=\frac{(\alpha)_2\cdot(\beta)_2\cdot(\gamma)_2}{(\delta)_2 \cdot (\varepsilon)_2}
        = \frac{\alpha(\alpha+1)\beta(\beta+1)\gamma(\gamma+1)}{\delta(\delta+1)\varepsilon(\varepsilon+1)}\\
        &= \frac{2a(2a+1)2b(2b+1)(a+b)(a+b+1)}{(2a+2b)(2a+2b+1)\left(a+b+\frac 1 2\right)\left(a+b+\frac 3 2\right)}\\
        &= \frac{ab(2a+1)(2b+1)(a+b+1)}{(a+b+1/2)^2(a+b+3/2)}\\
        &=\frac{ab(4a^2b+4ab^2+8ab+2a^2+2b^2+3a+3b+1)}{(a+b+1/2)^2(a+b+3/2)}=h''(0)
    \end{align*}
    \en{From the Picard-Lindelöf theorem we deduce that $h(z)=g(z)$ and thus we have proven Clausen's formula.}%
    \de{Insgesamt gilt also wegen des Satzes von Picard-Lindelöf $h(z)=g(z)$ und wir haben die Formel von Clausen bewiesen.}
\end{proof}


